
\noindent
Wir betrachten im Folgenden die Funktion
\begin{equation}
\label{eq:f_definition}
f(x) = \cos(2x^3 + 1)\,e^{x-1}
\end{equation}

\noindent
Eine Funktion $f$ lässt sich allgemein als Summe zweier Teilfunktionen $g$ und $h$
schreiben:
\begin{equation}
\label{eq:f_zerlegung}
f(x) = g(x) + h(x)
\end{equation}

\noindent
Dabei bezeichnet $g(x)$ den geraden Anteil und $h(x)$ den ungeraden Anteil der
Funktion. Diese Anteile sind definiert durch:
\begin{equation}
\label{eq:g_definition}
g(x) = \frac{f(x) + f(-x)}{2}
\end{equation}
\begin{equation}
\label{eq:h_definition}
h(x) = \frac{f(x) - f(-x)}{2}
\end{equation}

\noindent
Zur Anwendung dieser Zerlegung auf die gegebene Funktion wird zunächst $f(-x)$
bestimmt. Durch Einsetzen von $-x$ in Gleichung~\eqref{eq:f_definition} ergibt sich:
\begin{equation}
\label{eq:f_minus_x}
f(-x) = \cos(-2x^3 + 1)\,e^{-x-1}
\end{equation}

\noindent
Setzt man die Ausdrücke für $f(x)$ und $f(-x)$ in die Definition des geraden Anteils
nach Gleichung~\eqref{eq:g_definition} ein, erhält man:
\begin{equation}
\label{eq:g_result}
g(x) = \frac{1}{2}\left(\cos(2x^3+1)\,e^{x-1} + \cos(-2x^3+1)\,e^{-x-1}\right)
\end{equation}

\noindent
Durch Ausklammern von $e^{-1}$ kann der gerade Anteil äquivalent geschrieben werden
als:
\begin{equation}
\label{eq:g_result_factored}
g(x) = \frac{e^{-1}}{2}\left(\cos(2x^3+1)\,e^{x} + \cos(-2x^3+1)\,e^{-x}\right)
\end{equation}

\noindent
Analog ergibt sich für den ungeraden Anteil durch Einsetzen in Gleichung~\eqref{eq:h_definition}:
\begin{equation}
\label{eq:h_result}
h(x) = \frac{1}{2}\left(\cos(2x^3+1)\,e^{x-1} - \cos(-2x^3+1)\,e^{-x-1}\right)
\end{equation}

\noindent
Auch hier kann $e^{-1}$ ausgeklammert werden, sodass folgt:
\begin{equation}
\label{eq:h_result_factored}
h(x) = \frac{e^{-1}}{2}\left(\cos(2x^3+1)\,e^{x} - \cos(-2x^3+1)\,e^{-x}\right)
\end{equation}

\noindent
Damit ist die Funktion $f(x)$ eindeutig in einen geraden Anteil $g(x)$ und einen
ungeraden Anteil $h(x)$ zerlegt.
